\setlength{\baselineskip}{20pt}
\chapter{结论与展望}
\label{cha:chap5}

\section{全文工作总结}
本文围绕复杂视觉场景下判别模型的鲁棒性与可部署性问题，系统研究了扩散思想在判别表征学习中的应用。全文主要完成了以下工作：
\begin{enumerate}
  \item 提出 DenoiseRep 方法，将级联嵌入层与逐层去噪过程统一建模，实现特征提取与去噪增强协同优化；
  \item 设计参数融合机制，将训练阶段去噪分支等价折叠到推理阶段主干映射，实现无额外推理开销的特征增强；
  \item 提出 DenoiseRep++ 增强框架，通过教师--学生蒸馏和多步去噪融合提升轻量去噪模块的表达能力；
  \item 在 ReID、分类、检测与分割任务上完成系统实验，验证方法在性能提升、跨任务泛化与效率保持方面的综合优势。
\end{enumerate}

\section{主要创新点}
本文的主要创新可归纳为以下三个方面：
\begin{enumerate}
  \item \textbf{统一建模创新}：将扩散去噪机制引入判别表征学习主干，构建了生成思想与判别学习之间的统一桥接框架；
  \item \textbf{高效部署创新}：提出去噪层与特征层参数融合策略，在理论等价前提下实现训练增强与推理高效的兼容；
  \item \textbf{能力增强创新}：提出蒸馏与多步融合协同机制，提升了轻量结构在复杂噪声与分布偏移条件下的鲁棒性。
\end{enumerate}

\section{不足与未来工作}
尽管本文方法取得了较稳定的实验增益，但仍存在进一步优化空间：
\begin{enumerate}
  \item 当前参数融合主要基于线性近似，后续可探索更灵活的非线性融合机制；
  \item 多步融合策略在极低算力设备上的超参敏感性仍需进一步分析；
  \item 现阶段实验以视觉单模态为主，尚未系统扩展至图文联合或时序多模态任务。
\end{enumerate}

未来工作将重点围绕以下方向展开：
\begin{enumerate}
  \item 面向多模态表征学习扩展扩散去噪增强框架；
  \item 引入自适应噪声调度与动态融合策略，提升不同场景下的自适应能力；
  \item 面向真实部署场景建立更完整的鲁棒评测基准，进一步验证方法的工程可用性。
\end{enumerate}

总体而言，本文工作为扩散模型在判别任务中的高效应用提供了可复用的研究范式，并为复杂场景下高可靠视觉感知系统的构建提供了方法基础。
